% ============================================================================
% BỔ SUNG NGÂN HÀNG TOÁN 9 KNTT — CHƯƠNG V: ĐƯỜNG TRÒN (BÀI 14, 15, 16, 17)
% 11 mục nhỏ (dạng) × 3 câu (1 TN Nhận biết + 1 Đúng/Sai Thông hiểu + 1 TLN Vận dụng) = 33 câu
% ============================================================================

% ==================== BÀI 14. CUNG VÀ DÂY CỦA MỘT ĐƯỜNG TRÒN ====================

% --- 9H5?14-1: Quan hệ giữa đường kính và dây ---

\begin{ex}%[9H5N14-1]
% nguon: SGK Toán 9 KNTT tập 1 – Bài 14, Cung và dây của một đường tròn
	Cho đường tròn $(O; 6\text{ cm})$. Độ dài dây lớn nhất của đường tròn đó bằng
	\choice
	{$3\text{ cm}$}
	{$6\text{ cm}$}
	{\True $12\text{ cm}$}
	{$18\text{ cm}$}
	\loigiai{
		Trong một đường tròn, đường kính là dây lớn nhất. Do bán kính $R = 6\text{ cm}$ nên độ dài dây lớn nhất bằng $2R = 12\text{ cm}$.
	}
\end{ex}

\begin{ex}%[9H5H14-1]
% nguon: SGK Toán 9 KNTT tập 1 – Bài 14, Cung và dây của một đường tròn
	Cho đường tròn $(O; 5\text{ cm})$ có đường kính $AB$ và một dây $CD$ không đi qua tâm $O$ có độ dài $CD = 8\text{ cm}$. Xét tính đúng, sai của các khẳng định sau:
	\choiceTF
	{\True Độ dài đường kính $AB$ bằng $10\text{ cm}$}
	{Độ dài dây $CD$ lớn hơn độ dài đường kính $AB$}
	{\True Khoảng cách từ tâm $O$ đến dây $CD$ bằng $3\text{ cm}$}
	{Nếu một dây $EF$ của đường tròn $(O; 5\text{ cm})$ có độ dài $6\text{ cm}$ thì diện tích tam giác $OEF$ lớn hơn diện tích tam giác $OCD$}
	\loigiai{
		\begin{itemchoice}
			\itemch Đúng. Đường kính $AB = 2R = 2 \cdot 5 = 10\text{ cm}$.
			\itemch Sai. Trong một đường tròn, mọi dây không đi qua tâm đều ngắn hơn đường kính, cụ thể $CD = 8\text{ cm} < AB = 10\text{ cm}$.
			\itemch Đúng. Gọi $H$ là trung điểm của $CD$, ta có $OH \perp CD$ và $CH = \dfrac{CD}{2} = 4\text{ cm}$. Áp dụng định lí Pythagore cho tam giác $OHC$ vuông tại $H$: $OH = \sqrt{OC^2 - CH^2} = \sqrt{5^2 - 4^2} = 3\text{ cm}$.
			\itemch Sai. Khoảng cách từ $O$ đến dây $EF$ là $\sqrt{5^2 - 3^2} = 4\text{ cm}$. Ta có $S_{OCD} = \dfrac{1}{2} \cdot 8 \cdot 3 = 12\text{ cm}^2$ và $S_{OEF} = \dfrac{1}{2} \cdot 6 \cdot 4 = 12\text{ cm}^2$, hai diện tích bằng nhau.
		\end{itemchoice}
	}
\end{ex}

\begin{bt}%[9H5V14-1]
% nguon: SGK Toán 9 KNTT tập 1 – Bài 14, Cung và dây của một đường tròn
	Cho đường tròn tâm $O$ có đường kính $AB = 20\text{ cm}$. Dây $CD$ vuông góc với đường kính $AB$ tại điểm $H$ sao cho $AH = 4\text{ cm}$. Tính độ dài dây $CD$ (đơn vị cm).
	\shortans{16}
	\loigiai{
		Bán kính đường tròn là $R = \dfrac{AB}{2} = 10\text{ cm}$.\\
		Vì $H$ nằm giữa $A$ và $O$ nên $OH = OA - AH = 10 - 4 = 6\text{ cm}$.\\
		Do $AB \perp CD$ tại $H$ nên $H$ là trung điểm của $CD$. Xét tam giác $OHC$ vuông tại $H$:
		$$CH = \sqrt{OC^2 - OH^2} = \sqrt{10^2 - 6^2} = 8\text{ cm}.$$
		Vậy độ dài dây $CD = 2CH = 16\text{ cm}$.
	}
\end{bt}

% --- 9H5?14-2: Góc ở tâm, số đo cung ---

\begin{ex}%[9H5N14-2]
% nguon: SGK Toán 9 KNTT tập 1 – Bài 14, Cung và dây của một đường tròn
	Trên đường tròn $(O)$, lấy hai điểm $A, B$ sao cho $\widehat{AOB} = 75^\circ$. Số đo cung lớn $\overset{\frown}{AB}$ bằng
	\choice
	{$75^\circ$}
	{$105^\circ$}
	{$150^\circ$}
	{\True $285^\circ$}
	\loigiai{
		Số đo cung nhỏ $\overset{\frown}{AB}$ bằng số đo góc ở tâm chắn cung đó là $75^\circ$.\\
		Suy ra số đo cung lớn $\overset{\frown}{AB}$ bằng $360^\circ - 75^\circ = 285^\circ$.
	}
\end{ex}

\begin{ex}%[9H5H14-2]
% nguon: SGK Toán 9 KNTT tập 1 – Bài 14, Cung và dây của một đường tròn
	Trên đường tròn $(O; 4\text{ cm})$, lấy hai điểm $A, B$ sao cho $\widehat{AOB} = 120^\circ$ và lấy điểm $C$ nằm trên cung nhỏ $\overset{\frown}{AB}$ sao cho số đo cung nhỏ $\overset{\frown}{AC}$ bằng $50^\circ$. Xét tính đúng, sai của các khẳng định sau:
	\choiceTF
	{\True Số đo cung nhỏ $\overset{\frown}{AB}$ bằng $120^\circ$}
	{Số đo cung lớn $\overset{\frown}{AB}$ bằng $220^\circ$}
	{\True Số đo cung nhỏ $\overset{\frown}{BC}$ bằng $70^\circ$}
	{\True Độ dài dây $AB$ bằng $4\sqrt{3}\text{ cm}$}
	\loigiai{
		\begin{itemchoice}
			\itemch Đúng. Số đo cung nhỏ $\overset{\frown}{AB}$ bằng $\widehat{AOB} = 120^\circ$.
			\itemch Sai. Số đo cung lớn $\overset{\frown}{AB}$ bằng $360^\circ - 120^\circ = 240^\circ$.
			\itemch Đúng. Vì $C$ nằm trên cung nhỏ $\overset{\frown}{AB}$ nên số đo cung nhỏ $\overset{\frown}{BC}$ bằng $120^\circ - 50^\circ = 70^\circ$.
			\itemch Đúng. Kẻ $OH \perp AB$ tại $H$, tam giác $OAB$ cân tại $O$ có $\widehat{AOH} = 60^\circ$ nên $AH = OA \cdot \sin 60^\circ = 4 \cdot \dfrac{\sqrt{3}}{2} = 2\sqrt{3}\text{ cm}$, suy ra $AB = 2AH = 4\sqrt{3}\text{ cm}$.
		\end{itemchoice}
	}
\end{ex}

\begin{bt}%[9H5V14-2]
% nguon: SGK Toán 9 KNTT tập 1 – Bài 14, Cung và dây của một đường tròn
	Lúc $4$ giờ đúng, kim giờ và kim phút của một đồng hồ tạo thành góc ở tâm chắn cung nhỏ trên mặt đồng hồ có số đo bằng bao nhiêu độ?
	\shortans{120}
	\loigiai{
		Mặt đồng hồ chia đường tròn thành $12$ cung bằng nhau, mỗi cung tương ứng với $1$ giờ có số đo là $\dfrac{360^\circ}{12} = 30^\circ$.\\
		Lúc $4$ giờ đúng, kim phút chỉ số $12$ và kim giờ chỉ số $4$, khoảng cách giữa hai kim là $4$ khoảng giờ nên số đo cung nhỏ bị chắn là $4 \cdot 30^\circ = 120^\circ$.
	}
\end{bt}

% --- 9H5?14-3: Liên hệ giữa dây và khoảng cách đến tâm ---

\begin{ex}%[9H5N14-3]
% nguon: SGK Toán 9 KNTT tập 1 – Bài 14, Cung và dây của một đường tròn
	Cho đường tròn $(O; 13\text{ cm})$ và dây $AB = 24\text{ cm}$. Khoảng cách từ tâm $O$ đến dây $AB$ bằng
	\choice
	{\True $5\text{ cm}$}
	{$6\text{ cm}$}
	{$11\text{ cm}$}
	{$12\text{ cm}$}
	\loigiai{
		Gọi $H$ là trung điểm của dây $AB$, ta có $OH \perp AB$ và $AH = \dfrac{AB}{2} = 12\text{ cm}$.\\
		Trong tam giác $OHA$ vuông tại $H$: $OH = \sqrt{OA^2 - AH^2} = \sqrt{13^2 - 12^2} = 5\text{ cm}$.
	}
\end{ex}

\begin{ex}%[9H5H14-3]
% nguon: SGK Toán 9 KNTT tập 1 – Bài 14, Cung và dây của một đường tròn
	Cho đường tròn $(O; 10\text{ cm})$ có hai dây song song $AB = 12\text{ cm}$ và $CD = 16\text{ cm}$ nằm cùng một phía đối với tâm $O$. Xét tính đúng, sai của các khẳng định sau:
	\choiceTF
	{\True Khoảng cách từ tâm $O$ đến dây $AB$ bằng $8\text{ cm}$}
	{\True Khoảng cách từ tâm $O$ đến dây $CD$ bằng $6\text{ cm}$}
	{Khoảng cách giữa hai dây song song $AB$ và $CD$ bằng $14\text{ cm}$}
	{\True Diện tích hình thang $ABCD$ bằng $28\text{ cm}^2$}
	\loigiai{
		\begin{itemchoice}
			\itemch Đúng. Gọi $H$ là trung điểm $AB$, ta có $AH = 6\text{ cm}$ nên $OH = \sqrt{10^2 - 6^2} = 8\text{ cm}$.
			\itemch Đúng. Gọi $K$ là trung điểm $CD$, ta có $CK = 8\text{ cm}$ nên $OK = \sqrt{10^2 - 8^2} = 6\text{ cm}$.
			\itemch Sai. Vì hai dây $AB$ và $CD$ nằm cùng một phía đối với tâm $O$ nên khoảng cách giữa chúng là $HK = OH - OK = 8 - 6 = 2\text{ cm}$.
			\itemch Đúng. Hình thang $ABCD$ có hai đáy $AB = 12\text{ cm}$, $CD = 16\text{ cm}$ và chiều cao $HK = 2\text{ cm}$ nên diện tích là $\dfrac{12 + 16}{2} \cdot 2 = 28\text{ cm}^2$.
		\end{itemchoice}
	}
\end{ex}

\begin{bt}%[9H5V14-3]
% nguon: SGK Toán 9 KNTT tập 1 – Bài 14, Cung và dây của một đường tròn
	Một đường ống nước hình trụ nằm ngang có mặt cắt ngang là hình tròn bán kính $25\text{ cm}$. Khi mực nước trong ống sâu $10\text{ cm}$ (tính từ điểm thấp nhất của đáy ống lên mặt nước) thì bề rộng mặt thoáng của nước trong ống bằng bao nhiêu xăng-ti-mét?
	\shortans{40}
	\loigiai{
		Gọi $O$ là tâm mặt cắt tròn bán kính $R = 25\text{ cm}$, $AB$ là dây biểu thị mặt thoáng của nước.\\
		Khoảng cách từ tâm $O$ đến mặt thoáng $AB$ là $d = 25 - 10 = 15\text{ cm}$.\\
		Nửa bề rộng mặt thoáng là $\sqrt{25^2 - 15^2} = \sqrt{400} = 20\text{ cm}$.\\
		Vậy bề rộng mặt thoáng của nước là $AB = 2 \cdot 20 = 40\text{ cm}$.
	}
\end{bt}

% ==================== BÀI 15. ĐỘ DÀI CUNG TRÒN. DIỆN TÍCH QUẠT TRÒN, VÀNH KHUYÊN ====================

% --- 9H5?15-1: Độ dài đường tròn, cung tròn ---

\begin{ex}%[9H5N15-1]
% nguon: SGK Toán 9 KNTT tập 1 – Bài 15, Độ dài của cung tròn. Diện tích hình quạt tròn và hình vành khuyên
	Độ dài cung tròn $60^\circ$ của đường tròn bán kính $R = 9\text{ cm}$ bằng
	\choice
	{$\dfrac{3\pi}{2}\text{ cm}$}
	{\True $3\pi\text{ cm}$}
	{$6\pi\text{ cm}$}
	{$9\pi\text{ cm}$}
	\loigiai{
		Độ dài cung tròn $60^\circ$ bán kính $R = 9\text{ cm}$ là $l = \dfrac{n}{180}\pi R = \dfrac{60}{180}\pi \cdot 9 = 3\pi\text{ cm}$.
	}
\end{ex}

\begin{ex}%[9H5H15-1]
% nguon: SGK Toán 9 KNTT tập 1 – Bài 15, Độ dài của cung tròn. Diện tích hình quạt tròn và hình vành khuyên
	Cho đường tròn $(O; 12\text{ cm})$ và cung nhỏ $\overset{\frown}{AB}$ có số đo $90^\circ$. Xét tính đúng, sai của các khẳng định sau:
	\choiceTF
	{\True Chu vi của đường tròn $(O)$ bằng $24\pi\text{ cm}$}
	{\True Độ dài cung nhỏ $\overset{\frown}{AB}$ bằng $6\pi\text{ cm}$}
	{Độ dài cung lớn $\overset{\frown}{AB}$ bằng $16\pi\text{ cm}$}
	{\True Chu vi của hình quạt tròn $OAB$ (ứng với cung nhỏ $\overset{\frown}{AB}$), lấy $\pi \approx 3{,}14$ và làm tròn đến hàng phần mười, bằng $42{,}8\text{ cm}$}
	\loigiai{
		\begin{itemchoice}
			\itemch Đúng. Chu vi đường tròn là $C = 2\pi R = 2\pi \cdot 12 = 24\pi\text{ cm}$.
			\itemch Đúng. Độ dài cung nhỏ $\overset{\frown}{AB}$ là $l = \dfrac{90}{180}\pi \cdot 12 = 6\pi\text{ cm}$.
			\itemch Sai. Độ dài cung lớn $\overset{\frown}{AB}$ là $24\pi - 6\pi = 18\pi\text{ cm}$.
			\itemch Đúng. Chu vi hình quạt tròn $OAB$ bằng $2R + l = 2 \cdot 12 + 6\pi \approx 24 + 6 \cdot 3{,}14 = 42{,}84 \approx 42{,}8\text{ cm}$.
		\end{itemchoice}
	}
\end{ex}

\begin{bt}%[9H5V15-1]
% nguon: SGK Toán 9 KNTT tập 1 – Bài 15, Độ dài của cung tròn. Diện tích hình quạt tròn và hình vành khuyên
	Một bánh xe đạp có bán kính $35\text{ cm}$. Khi bánh xe lăn không trượt được $10$ vòng trên một đoạn đường thẳng thì xe đi được quãng đường dài bao nhiêu mét? (Lấy $\pi \approx 3{,}14$, làm tròn kết quả đến hàng đơn vị).
	\shortans{22}
	\loigiai{
		Đổi $R = 35\text{ cm} = 0{,}35\text{ m}$.\\
		Chu vi của bánh xe là $C = 2\pi R \approx 2 \cdot 3{,}14 \cdot 0{,}35 = 2{,}198\text{ m}$.\\
		Quãng đường xe đi được sau $10$ vòng lăn là $10 \cdot 2{,}198 = 21{,}98 \approx 22\text{ m}$.
	}
\end{bt}

% --- 9H5?15-2: Diện tích hình tròn, hình quạt tròn ---

\begin{ex}%[9H5N15-2]
% nguon: SGK Toán 9 KNTT tập 1 – Bài 15, Độ dài của cung tròn. Diện tích hình quạt tròn và hình vành khuyên
	Diện tích hình quạt tròn bán kính $R = 6\text{ cm}$, ứng với cung có số đo $120^\circ$ bằng
	\choice
	{$4\pi\text{ cm}^2$}
	{$6\pi\text{ cm}^2$}
	{\True $12\pi\text{ cm}^2$}
	{$24\pi\text{ cm}^2$}
	\loigiai{
		Diện tích hình quạt tròn là $S = \dfrac{n}{360}\pi R^2 = \dfrac{120}{360}\pi \cdot 6^2 = 12\pi\text{ cm}^2$.
	}
\end{ex}

\begin{ex}%[9H5H15-2]
% nguon: SGK Toán 9 KNTT tập 1 – Bài 15, Độ dài của cung tròn. Diện tích hình quạt tròn và hình vành khuyên
	Cho hình quạt tròn tâm $O$, bán kính $R = 6\text{ cm}$, ứng với cung nhỏ $\overset{\frown}{AB}$ có số đo $60^\circ$. Xét tính đúng, sai của các khẳng định sau:
	\choiceTF
	{\True Diện tích hình tròn $(O; 6\text{ cm})$ bằng $36\pi\text{ cm}^2$}
	{\True Diện tích hình quạt tròn $OAB$ bằng $6\pi\text{ cm}^2$}
	{Diện tích tam giác $OAB$ bằng $18\sqrt{3}\text{ cm}^2$}
	{\True Diện tích hình viên phân giới hạn bởi dây $AB$ và cung nhỏ $\overset{\frown}{AB}$ (lấy $\pi \approx 3{,}14$, $\sqrt{3} \approx 1{,}73$ và làm tròn đến hàng phần mười) bằng $3{,}3\text{ cm}^2$}
	\loigiai{
		\begin{itemchoice}
			\itemch Đúng. Diện tích hình tròn là $\pi R^2 = \pi \cdot 6^2 = 36\pi\text{ cm}^2$.
			\itemch Đúng. Diện tích hình quạt tròn là $\dfrac{60}{360}\pi \cdot 6^2 = 6\pi\text{ cm}^2$.
			\itemch Sai. Vì $OA = OB = 6\text{ cm}$ và $\widehat{AOB} = 60^\circ$ nên tam giác $OAB$ đều cạnh $6\text{ cm}$, có diện tích $\dfrac{6^2\sqrt{3}}{4} = 9\sqrt{3}\text{ cm}^2$.
			\itemch Đúng. Diện tích hình viên phân bằng $6\pi - 9\sqrt{3} \approx 6 \cdot 3{,}14 - 9 \cdot 1{,}73 = 18{,}84 - 15{,}57 = 3{,}27 \approx 3{,}3\text{ cm}^2$.
		\end{itemchoice}
	}
\end{ex}

\begin{bt}%[9H5V15-2]
% nguon: SGK Toán 9 KNTT tập 1 – Bài 15, Độ dài của cung tròn. Diện tích hình quạt tròn và hình vành khuyên
	Một vòi phun nước tự động quay đều qua một góc $150^\circ$ và có bán kính phun xa tối đa $6\text{ m}$. Tính diện tích phần bãi cỏ được vòi nước tưới (đơn vị $\text{m}^2$, lấy $\pi \approx 3{,}14$, làm tròn kết quả đến hàng phần mười).
	\shortans{47,1}
	\loigiai{
		Phần bãi cỏ được tưới là hình quạt tròn bán kính $R = 6\text{ m}$, góc ở tâm $n = 150^\circ$.\\
		Diện tích được tưới là $S = \dfrac{150}{360}\pi \cdot 6^2 = 15\pi \approx 15 \cdot 3{,}14 = 47{,}1\text{ m}^2$.
	}
\end{bt}

% --- 9H5?15-3: Diện tích hình vành khuyên, bài toán thực tế ---

\begin{ex}%[9H5N15-3]
% nguon: SGK Toán 9 KNTT tập 1 – Bài 15, Độ dài của cung tròn. Diện tích hình quạt tròn và hình vành khuyên
	Diện tích hình vành khuyên giới hạn bởi hai đường tròn đồng tâm có bán kính $R = 5\text{ cm}$ và $r = 3\text{ cm}$ bằng
	\choice
	{$2\pi\text{ cm}^2$}
	{$4\pi\text{ cm}^2$}
	{\True $16\pi\text{ cm}^2$}
	{$34\pi\text{ cm}^2$}
	\loigiai{
		Diện tích hình vành khuyên là $S = \pi(R^2 - r^2) = \pi(5^2 - 3^2) = 16\pi\text{ cm}^2$.
	}
\end{ex}

\begin{ex}%[9H5H15-3]
% nguon: SGK Toán 9 KNTT tập 1 – Bài 15, Độ dài của cung tròn. Diện tích hình quạt tròn và hình vành khuyên
	Một bồn hoa hình tròn có bán kính $4\text{ m}$. Xung quanh bồn hoa người ta xây một lối đi hình vành khuyên có bề rộng $1\text{ m}$. Xét tính đúng, sai của các khẳng định sau:
	\choiceTF
	{\True Bán kính đường tròn ngoài của lối đi bằng $5\text{ m}$}
	{\True Diện tích của bồn hoa bằng $16\pi\text{ m}^2$}
	{Diện tích của lối đi bằng $10\pi\text{ m}^2$}
	{\True Nếu chi phí lát gạch lối đi là $200\,000\text{ đồng/m}^2$ (lấy $\pi \approx 3{,}14$) thì tổng chi phí lát gạch lối đi là $5\,652\,000$ đồng}
	\loigiai{
		\begin{itemchoice}
			\itemch Đúng. Bán kính vòng tròn ngoài là $R = 4 + 1 = 5\text{ m}$.
			\itemch Đúng. Diện tích bồn hoa là $\pi r^2 = \pi \cdot 4^2 = 16\pi\text{ m}^2$.
			\itemch Sai. Diện tích lối đi hình vành khuyên là $\pi(R^2 - r^2) = \pi(5^2 - 4^2) = 9\pi\text{ m}^2$.
			\itemch Đúng. Tổng chi phí lát gạch là $9 \cdot 3{,}14 \cdot 200\,000 = 5\,652\,000$ đồng.
		\end{itemchoice}
	}
\end{ex}

\begin{bt}%[9H5V15-3]
% nguon: SGK Toán 9 KNTT tập 1 – Bài 15, Độ dài của cung tròn. Diện tích hình quạt tròn và hình vành khuyên
	Một chiếc đĩa CD có đường kính ngoài $12\text{ cm}$ và lỗ tròn đồng tâm ở giữa có đường kính $2\text{ cm}$. Tính diện tích một mặt đĩa (phần hình vành khuyên, đơn vị $\text{cm}^2$, lấy $\pi \approx 3{,}14$, làm tròn kết quả đến hàng đơn vị).
	\shortans{110}
	\loigiai{
		Bán kính ngoài của đĩa là $R = \dfrac{12}{2} = 6\text{ cm}$, bán kính lỗ tròn ở giữa là $r = \dfrac{2}{2} = 1\text{ cm}$.\\
		Diện tích một mặt đĩa là $S = \pi(R^2 - r^2) = \pi(6^2 - 1^2) = 35\pi \approx 35 \cdot 3{,}14 = 109{,}9 \approx 110\text{ cm}^2$.
	}
\end{bt}

% ==================== BÀI 16. VỊ TRÍ TƯƠNG ĐỐI CỦA ĐƯỜNG THẲNG VÀ ĐƯỜNG TRÒN ====================

% --- 9H5?16-1: Xác định vị trí tương đối của đường thẳng và đường tròn ---

\begin{ex}%[9H5N16-1]
% nguon: SGK Toán 9 KNTT tập 1 – Bài 16, Vị trí tương đối của đường thẳng và đường tròn
	Cho đường tròn $(O; 5\text{ cm})$ và đường thẳng $a$ có khoảng cách từ tâm $O$ đến $a$ bằng $7\text{ cm}$. Vị trí tương đối của đường thẳng $a$ và đường tròn $(O)$ là
	\choice
	{Cắt nhau}
	{Tiếp xúc nhau}
	{\True Không giao nhau}
	{Đi qua tâm $O$}
	\loigiai{
		Vì khoảng cách từ tâm $O$ đến đường thẳng $a$ là $d = 7\text{ cm} > R = 5\text{ cm}$ nên đường thẳng $a$ và đường tròn $(O)$ không giao nhau.
	}
\end{ex}

\begin{ex}%[9H5H16-1]
% nguon: SGK Toán 9 KNTT tập 1 – Bài 16, Vị trí tương đối của đường thẳng và đường tròn
	Cho đường tròn $(O; 6\text{ cm})$ và ba đường thẳng $a, b, c$ có khoảng cách đến tâm $O$ lần lượt là $d_a = 4\text{ cm}$, $d_b = 6\text{ cm}$, $d_c = 8\text{ cm}$. Xét tính đúng, sai của các khẳng định sau:
	\choiceTF
	{\True Đường thẳng $a$ cắt đường tròn $(O)$ tại hai điểm phân biệt}
	{\True Đường thẳng $b$ là tiếp tuyến của đường tròn $(O)$}
	{Đường thẳng $c$ và đường tròn $(O)$ có đúng một điểm chung}
	{\True Độ dài dây cung tạo bởi đường thẳng $a$ và đường tròn $(O)$ bằng $4\sqrt{5}\text{ cm}$}
	\loigiai{
		\begin{itemchoice}
			\itemch Đúng. Vì $d_a = 4\text{ cm} < R = 6\text{ cm}$ nên đường thẳng $a$ cắt đường tròn $(O)$ tại hai điểm phân biệt.
			\itemch Đúng. Vì $d_b = 6\text{ cm} = R$ nên đường thẳng $b$ tiếp xúc với đường tròn $(O)$, tức là tiếp tuyến của $(O)$.
			\itemch Sai. Vì $d_c = 8\text{ cm} > R = 6\text{ cm}$ nên đường thẳng $c$ và đường tròn $(O)$ không có điểm chung nào.
			\itemch Đúng. Độ dài dây cung cắt bởi đường thẳng $a$ là $2\sqrt{R^2 - d_a^2} = 2\sqrt{6^2 - 4^2} = 2\sqrt{20} = 4\sqrt{5}\text{ cm}$.
		\end{itemchoice}
	}
\end{ex}

\begin{bt}%[9H5V16-1]
% nguon: SGK Toán 9 KNTT tập 1 – Bài 16, Vị trí tương đối của đường thẳng và đường tròn
	Trong mặt phẳng tọa độ $Oxy$, cho điểm $A(3; 4)$. Gọi $R$ là bán kính của đường tròn tâm $A$ sao cho đường tròn đó tiếp xúc với trục tung $Oy$. Giá trị của $R$ bằng bao nhiêu?
	\shortans{3}
	\loigiai{
		Khoảng cách từ điểm $A(3; 4)$ đến trục tung $Oy$ bằng $|x_A| = |3| = 3$.\\
		Để đường tròn $(A; R)$ tiếp xúc với trục tung $Oy$ thì $R = d(A, Oy) = 3$.
	}
\end{bt}

% --- 9H5?16-2: Tiếp tuyến: dấu hiệu nhận biết, chứng minh tiếp tuyến ---

\begin{ex}%[9H5N16-2]
% nguon: SGK Toán 9 KNTT tập 1 – Bài 16, Vị trí tương đối của đường thẳng và đường tròn
	Đường thẳng $d$ là tiếp tuyến của đường tròn $(O; R)$ tại điểm $A$ khi và chỉ khi
	\choice
	{$A$ nằm trên đường tròn $(O)$}
	{$d$ đi qua tâm $O$}
	{\True $A$ nằm trên đường tròn $(O)$ và $d \perp OA$ tại $A$}
	{$OA = 2R$ và $d \perp OA$ tại $A$}
	\loigiai{
		Theo dấu hiệu nhận biết tiếp tuyến, đường thẳng $d$ là tiếp tuyến của đường tròn $(O; R)$ tại điểm $A$ khi và chỉ khi $A \in (O)$ và $d \perp OA$ tại $A$.
	}
\end{ex}

\begin{ex}%[9H5H16-2]
% nguon: SGK Toán 9 KNTT tập 1 – Bài 16, Vị trí tương đối của đường thẳng và đường tròn
	Cho tam giác $ABC$ có $AB = 6\text{ cm}$, $AC = 8\text{ cm}$, $BC = 10\text{ cm}$. Xét đường tròn tâm $C$ bán kính $R = 8\text{ cm}$. Xét tính đúng, sai của các khẳng định sau:
	\choiceTF
	{\True Tam giác $ABC$ vuông tại $A$}
	{\True Điểm $A$ nằm trên đường tròn $(C; 8\text{ cm})$}
	{Đường thẳng $AB$ cắt đường tròn $(C; 8\text{ cm})$ tại hai điểm phân biệt}
	{\True Khoảng cách từ điểm $A$ đến đường thẳng $BC$ bằng $4{,}8\text{ cm}$}
	\loigiai{
		\begin{itemchoice}
			\itemch Đúng. Ta có $AB^2 + AC^2 = 6^2 + 8^2 = 100 = 10^2 = BC^2$, theo định lí Pythagore đảo thì tam giác $ABC$ vuông tại $A$.
			\itemch Đúng. Vì $CA = 8\text{ cm} = R$ nên điểm $A$ thuộc đường tròn $(C; 8\text{ cm})$.
			\itemch Sai. Vì $A \in (C; 8\text{ cm})$ và $AB \perp CA$ tại $A$ nên đường thẳng $AB$ là tiếp tuyến của đường tròn $(C; 8\text{ cm})$, chỉ có một điểm chung là $A$.
			\itemch Đúng. Đường cao $AH$ của tam giác vuông $ABC$ thoả mãn $AH \cdot BC = AB \cdot AC \Rightarrow AH = \dfrac{6 \cdot 8}{10} = 4{,}8\text{ cm}$.
		\end{itemchoice}
	}
\end{ex}

\begin{bt}%[9H5V16-2]
% nguon: SGK Toán 9 KNTT tập 1 – Bài 16, Vị trí tương đối của đường thẳng và đường tròn
	Cho đường tròn $(O; 5\text{ cm})$. Từ điểm $M$ nằm ngoài đường tròn sao cho $OM = 13\text{ cm}$, kẻ tiếp tuyến $MA$ đến đường tròn $(O)$ với $A$ là tiếp điểm. Tính độ dài đoạn thẳng $MA$ (đơn vị cm).
	\shortans{12}
	\loigiai{
		Vì $MA$ là tiếp tuyến của $(O)$ tại $A$ nên $MA \perp OA$, suy ra tam giác $OAM$ vuông tại $A$.\\
		Áp dụng định lí Pythagore: $MA = \sqrt{OM^2 - OA^2} = \sqrt{13^2 - 5^2} = \sqrt{144} = 12\text{ cm}$.
	}
\end{bt}

% --- 9H5?16-3: Tính chất hai tiếp tuyến cắt nhau ---

\begin{ex}%[9H5N16-3]
% nguon: SGK Toán 9 KNTT tập 1 – Bài 16, Vị trí tương đối của đường thẳng và đường tròn
	Từ điểm $M$ nằm ngoài đường tròn $(O)$, kẻ hai tiếp tuyến $MA, MB$ đến đường tròn ($A, B$ là các tiếp điểm). Khẳng định nào sau đây \textbf{sai}?
	\choice
	{$MA = MB$}
	{$MO$ là tia phân giác của $\widehat{AMB}$}
	{$OM$ là tia phân giác của $\widehat{AOB}$}
	{\True $OA = MA$}
	\loigiai{
		Theo tính chất hai tiếp tuyến cắt nhau, ta có $MA = MB$, $MO$ là tia phân giác của $\widehat{AMB}$ và $OM$ là tia phân giác của $\widehat{AOB}$. Độ dài $OA$ và $MA$ không nhất thiết bằng nhau (chỉ bằng nhau khi $\widehat{AMB} = 90^\circ$).
	}
\end{ex}

\begin{ex}%[9H5H16-3]
% nguon: SGK Toán 9 KNTT tập 1 – Bài 16, Vị trí tương đối của đường thẳng và đường tròn
	Từ điểm $M$ nằm ngoài đường tròn $(O; 3\text{ cm})$ sao cho $OM = 6\text{ cm}$, kẻ hai tiếp tuyến $MA, MB$ đến đường tròn ($A, B$ là các tiếp điểm). Xét tính đúng, sai của các khẳng định sau:
	\choiceTF
	{\True $MA = MB = 3\sqrt{3}\text{ cm}$}
	{\True $\widehat{AMO} = 30^\circ$}
	{Tam giác $MAB$ là tam giác vuông cân}
	{\True Chu vi tứ giác $OAMB$ bằng $(6 + 6\sqrt{3})\text{ cm}$}
	\loigiai{
		\begin{itemchoice}
			\itemch Đúng. Trong tam giác $OAM$ vuông tại $A$: $MA = \sqrt{OM^2 - OA^2} = \sqrt{6^2 - 3^2} = 3\sqrt{3}\text{ cm}$, và theo tính chất hai tiếp tuyến cắt nhau thì $MB = MA = 3\sqrt{3}\text{ cm}$.
			\itemch Đúng. Ta có $\sin \widehat{AMO} = \dfrac{OA}{OM} = \dfrac{3}{6} = \dfrac{1}{2}$ nên $\widehat{AMO} = 30^\circ$.
			\itemch Sai. Vì $MO$ là tia phân giác $\widehat{AMB}$ nên $\widehat{AMB} = 2 \cdot 30^\circ = 60^\circ$, kết hợp $MA = MB$ suy ra tam giác $MAB$ là tam giác đều.
			\itemch Đúng. Chu vi tứ giác $OAMB$ là $OA + OB + MA + MB = 3 + 3 + 3\sqrt{3} + 3\sqrt{3} = 6 + 6\sqrt{3}\text{ cm}$.
		\end{itemchoice}
	}
\end{ex}

\begin{bt}%[9H5V16-3]
% nguon: SGK Toán 9 KNTT tập 1 – Bài 16, Vị trí tương đối của đường thẳng và đường tròn
	Từ điểm $M$ nằm ngoài đường tròn $(O; 6\text{ cm})$ sao cho $OM = 10\text{ cm}$, kẻ hai tiếp tuyến $MA, MB$ đến đường tròn $(O)$ với $A, B$ là các tiếp điểm. Tính độ dài dây $AB$ (đơn vị cm, viết kết quả dưới dạng số thập phân).
	\shortans{9,6}
	\loigiai{
		Trong tam giác $OAM$ vuông tại $A$: $MA = \sqrt{OM^2 - OA^2} = \sqrt{10^2 - 6^2} = 8\text{ cm}$.\\
		Vì $OA = OB$ và $MA = MB$ nên $OM$ là đường trung trực của đoạn $AB$. Gọi $H$ là giao điểm của $OM$ và $AB$, ta có $AH \perp OM$ và $AB = 2AH$.\\
		Trong tam giác $OAM$ vuông tại $A$ có đường cao $AH$:
		$$AH = \dfrac{OA \cdot MA}{OM} = \dfrac{6 \cdot 8}{10} = 4{,}8\text{ cm}.$$
		Suy ra $AB = 2 \cdot 4{,}8 = 9{,}6\text{ cm}$.
	}
\end{bt}

% ==================== BÀI 17. VỊ TRÍ TƯƠNG ĐỐI CỦA HAI ĐƯỜNG TRÒN ====================

% --- 9H5?17-1: Xác định vị trí tương đối của hai đường tròn ---

\begin{ex}%[9H5N17-1]
% nguon: SGK Toán 9 KNTT tập 1 – Bài 17, Vị trí tương đối của hai đường tròn
	Cho hai đường tròn $(O; 7\text{ cm})$ và $(O'; 3\text{ cm})$ có khoảng cách hai tâm $OO' = 10\text{ cm}$. Vị trí tương đối của hai đường tròn đó là
	\choice
	{Ở ngoài nhau}
	{\True Tiếp xúc ngoài}
	{Cắt nhau}
	{Tiếp xúc trong}
	\loigiai{
		Ta có $R + r = 7 + 3 = 10\text{ cm} = OO'$, nên hai đường tròn $(O)$ và $(O')$ tiếp xúc ngoài với nhau.
	}
\end{ex}

\begin{ex}%[9H5H17-1]
% nguon: SGK Toán 9 KNTT tập 1 – Bài 17, Vị trí tương đối của hai đường tròn
	Cho hai đường tròn $(O; 8\text{ cm})$ và $(O'; 5\text{ cm})$ có khoảng cách giữa hai tâm là $OO' = d$. Xét tính đúng, sai của các khẳng định sau:
	\choiceTF
	{\True Nếu $d = 15\text{ cm}$ thì hai đường tròn ở ngoài nhau}
	{\True Nếu $d = 3\text{ cm}$ thì hai đường tròn tiếp xúc trong}
	{Nếu $d = 2\text{ cm}$ thì hai đường tròn cắt nhau}
	{\True Hai đường tròn cắt nhau khi và chỉ khi $3 < d < 13$}
	\loigiai{
		\begin{itemchoice}
			\itemch Đúng. Ta có $R + r = 8 + 5 = 13\text{ cm}$. Vì $d = 15 > 13$ nên hai đường tròn ở ngoài nhau.
			\itemch Đúng. Ta có $R - r = 8 - 5 = 3\text{ cm}$. Vì $d = 3 = R - r$ nên hai đường tròn tiếp xúc trong.
			\itemch Sai. Vì $d = 2 < R - r = 3$ nên hai đường tròn đựng nhau (không giao nhau).
			\itemch Đúng. Hai đường tròn cắt nhau khi và chỉ khi $R - r < d < R + r$, tức là $3 < d < 13$.
		\end{itemchoice}
	}
\end{ex}

\begin{bt}%[9H5V17-1]
% nguon: SGK Toán 9 KNTT tập 1 – Bài 17, Vị trí tương đối của hai đường tròn
	Cho hai đường tròn $(O; 10\text{ cm})$ và $(O'; 17\text{ cm})$ cắt nhau tại hai điểm $A, B$ sao cho dây chung $AB = 16\text{ cm}$ và hai tâm $O, O'$ nằm về hai phía của dây $AB$. Tính khoảng cách $OO'$ giữa hai tâm (đơn vị cm).
	\shortans{21}
	\loigiai{
		Vì $OO'$ là đường trung trực của dây chung $AB$ nên gọi $H = OO' \cap AB$, ta có $OO' \perp AB$ tại $H$ và $AH = \dfrac{AB}{2} = 8\text{ cm}$.\\
		Trong tam giác $OHA$ vuông tại $H$: $OH = \sqrt{OA^2 - AH^2} = \sqrt{10^2 - 8^2} = 6\text{ cm}$.\\
		Trong tam giác $O'HA$ vuông tại $H$: $O'H = \sqrt{O'A^2 - AH^2} = \sqrt{17^2 - 8^2} = 15\text{ cm}$.\\
		Vì $O$ và $O'$ nằm về hai phía của dây $AB$ nên $OO' = OH + O'H = 6 + 15 = 21\text{ cm}$.
	}
\end{bt}

% --- 9H5?17-2: Tiếp tuyến chung, bài toán tổng hợp ---

\begin{ex}%[9H5N17-2]
% nguon: SGK Toán 9 KNTT tập 1 – Bài 17, Vị trí tương đối của hai đường tròn
	Hai đường tròn ở ngoài nhau có tất cả bao nhiêu tiếp tuyến chung?
	\choice
	{$1$}
	{$2$}
	{$3$}
	{\True $4$}
	\loigiai{
		Hai đường tròn ở ngoài nhau có $2$ tiếp tuyến chung ngoài và $2$ tiếp tuyến chung trong, tổng cộng là $4$ tiếp tuyến chung.
	}
\end{ex}

\begin{ex}%[9H5H17-2]
% nguon: SGK Toán 9 KNTT tập 1 – Bài 17, Vị trí tương đối của hai đường tròn
	Cho hai đường tròn $(O; 9\text{ cm})$ và $(O'; 4\text{ cm})$ tiếp xúc ngoài tại điểm $A$. Gọi $BC$ là một tiếp tuyến chung ngoài của hai đường tròn ($B \in (O)$, $C \in (O')$). Xét tính đúng, sai của các khẳng định sau:
	\choiceTF
	{\True Độ dài đoạn nối tâm $OO'$ bằng $13\text{ cm}$}
	{\True Tứ giác $OBCO'$ là hình thang vuông}
	{Độ dài đoạn tiếp tuyến chung ngoài $BC$ bằng $10\text{ cm}$}
	{\True Diện tích tứ giác $OBCO'$ bằng $78\text{ cm}^2$}
	\loigiai{
		\begin{itemchoice}
			\itemch Đúng. Vì hai đường tròn tiếp xúc ngoài nên $OO' = R + r = 9 + 4 = 13\text{ cm}$.
			\itemch Đúng. Vì $OB \perp BC$ và $O'C \perp BC$ nên $OB \parallel O'C$ và $\widehat{OBC} = \widehat{O'CB} = 90^\circ$, do đó $OBCO'$ là hình thang vuông.
			\itemch Sai. Kẻ $O'H \perp OB$ tại $H$, khi đó $HBCO'$ là hình chữ nhật nên $BC = O'H$ và $OH = OB - O'C = 9 - 4 = 5\text{ cm}$. Áp dụng định lí Pythagore cho tam giác $OHO'$ vuông tại $H$: $BC = O'H = \sqrt{OO'^2 - OH^2} = \sqrt{13^2 - 5^2} = 12\text{ cm}$.
			\itemch Đúng. Diện tích hình thang vuông $OBCO'$ bằng $\dfrac{OB + O'C}{2} \cdot BC = \dfrac{9 + 4}{2} \cdot 12 = 78\text{ cm}^2$.
		\end{itemchoice}
	}
\end{ex}

\begin{bt}%[9H5V17-2]
% nguon: SGK Toán 9 KNTT tập 1 – Bài 17, Vị trí tương đối của hai đường tròn
	Hai bánh răng hình tròn có bán kính $R = 8\text{ cm}$ và $r = 2\text{ cm}$ đặt tiếp xúc ngoài nhau. Một đoạn dây truyền động căng thẳng tiếp xúc ngoài với cả hai bánh răng tại hai điểm $A$ và $B$. Tính độ dài đoạn thẳng $AB$ (đơn vị cm).
	\shortans{8}
	\loigiai{
		Gọi $O$ và $O'$ lần lượt là tâm hai bánh răng bán kính $R = 8\text{ cm}$ và $r = 2\text{ cm}$.\\
		Vì hai đường tròn tiếp xúc ngoài nên $OO' = R + r = 8 + 2 = 10\text{ cm}$.\\
		Kẻ $O'H \perp OA$ tại $H$, khi đó $HABO'$ là hình chữ nhật nên $AB = O'H$ và $OH = R - r = 8 - 2 = 6\text{ cm}$.\\
		Trong tam giác $OHO'$ vuông tại $H$: $AB = O'H = \sqrt{OO'^2 - OH^2} = \sqrt{10^2 - 6^2} = 8\text{ cm}$.
	}
\end{bt}
